\section{Related Work}

\para{Long-context handling for coding LLMs.}
Repository-level coding benchmarks such as SWE-bench cast software engineering as generating a patch for a real GitHub issue given a full codebase snapshot and an execution-based test oracle \cite{jimenez2023swebench,openai2024swebenchverified}. Agentic systems typically mitigate the resulting long-context burden through three complementary strategies. First, \emph{retrieval}: systems such as SWE-agent \cite{yang2024sweagent}, Agentless \cite{xia2024agentless}, and AutoCodeRover \cite{zhang2024autocoderover} search for relevant files, functions, or code snippets and present only these to the model, avoiding the need to pass the entire repository as context. Second, \emph{pruning and selection}: agents discard low-utility lines, constrain editable regions, or truncate long tool outputs to fit within context budgets \cite{yao2022react,schick2023toolformer}. Third, \emph{summarization and compression}: methods like LLMLingua replace long observations with model-generated summaries or compressed prompts to reduce token counts \cite{jiang2023llmlingua}. Retrieval-augmented generation is also a standard paradigm for selecting supporting context under a budget \cite{lewis2020rag,karpukhin2020dpr,robertson2009bm25}, and in code-centric settings, repository-level retrieval and representation learning are commonly used to identify relevant files or snippets \cite{zhang2023repocoder,feng2020codebert,guo2020graphcodebert,guo2022unixcoder}. Orthogonally, recent work on extending context windows through positional encoding modifications (e.g., RoPE scaling \cite{su2024rope}) and efficient attention mechanisms \cite{dao2022flashattention} has pushed context limits to 128K+ tokens, but these approaches do not reduce the per-token cost of processing long inputs. Our work explores a different axis: changing the \emph{modality} of long code context from text to images.

\para{Optical compression via image encoding.}
A separate line of work explores representing long textual contexts through the visual channel: text is rendered into a 2D layout and consumed as an image so that a vision-capable model can decode or reason over it with a smaller (or differently priced) input representation \cite{openai2023gpt4v_systemcard,shi2023gpt4v_ocr,kim2021donut,lee2022pix2struct}. DeepSeek-OCR explicitly frames this as \emph{contexts optical compression}, training a specialized encoder to keep the number of vision tokens manageable under high-resolution inputs while retaining decodability \cite{wei2025deepseekocr}. AgentOCR applies optical representations to compress agent interaction histories, showing that rendering prior tool outputs as images can reduce context length without significantly degrading task performance \cite{feng2026agentocr}. In the document understanding domain, models such as Nougat \cite{blecher2023nougat} and Pix2Struct \cite{lee2022pix2struct} are trained end-to-end to parse rendered document images, achieving strong results on scientific papers and web pages. Importantly, most reported gains in this space rely on specialized training or tailored visual encoders, leaving open how much benefit can be obtained \emph{without} post-training when applying optical encoding to code-centric artifacts \cite{wei2025deepseekocr,feng2026agentocr}.

\para{Vision-capable LLMs for code tasks.}
Vision-language models such as GPT-4V \cite{openai2023gpt4v_systemcard} and its successors have demonstrated the ability to read and reason about code from screenshots, UI mockups, and handwritten notes. Several studies have benchmarked these models on code-from-image tasks: for example, generating code from UI screenshots \cite{si2024design2code} or understanding code structure from rendered images \cite{shi2023gpt4v_ocr}. However, these evaluations typically focus on the model's vision capability as a \emph{standalone skill}---asking the model to read and reproduce code from an image in a single turn. Our work differs in that we use vision as an \emph{input channel within a multi-step agentic workflow}: the model must not only read code from images but also reason about it, plan edits, and generate patches over many turns of interaction. This places different demands on the vision encoder, requiring sustained fidelity across dozens of images per episode rather than one-shot transcription accuracy.

\para{Agentic coding and observation formatting.}
The design of the agent-computer interface (ACI) has been shown to significantly affect agentic coding performance. SWE-agent \cite{yang2024sweagent} introduced the concept of ACI design, demonstrating that the format in which tool outputs are presented to the model---including truncation, line numbering, and context windowing---can substantially change resolve rates. The ReAct framework \cite{yao2022react} established the interleaved reasoning-and-acting paradigm that most coding agents follow, while Reflexion \cite{shinn2023reflexion} showed that feeding back execution observations in structured formats improves self-correction. Our work contributes to this line by demonstrating that the \emph{modality} of observation formatting (text vs.\ image) is another axis of ACI design, and that changing this modality can alter model behavior in unexpected ways (Section~\ref{sec:prompt-alignment}), even when the underlying information content is preserved.

\para{Multimodal inputs and model behavior.}
An emerging body of work studies how the presence of multimodal inputs affects LLM behavior beyond the content of those inputs. Studies have shown that including images in prompts can alter the model's tone, verbosity, and instruction-following behavior, even when the images are irrelevant to the task \cite{gong2023figstep}. Our finding that image inputs cause GPT-5-mini to systematically change its git workflow (committing before diffing, Section~\ref{sec:prompt-alignment}) is consistent with this phenomenon: the presence of visual content appears to shift the model's behavioral distribution in ways that are independent of the image content itself. This suggests that multimodal agentic systems require careful prompt alignment to control for modality-induced behavioral confounds, a consideration that has received limited attention in the agentic coding literature.

\para{How our study differs.}
Rather than proposing a new trained compressor, we provide a controlled evaluation of optical encoding as a \emph{pure input representation change} for agentic coding. We test a fixed vision-capable model (GPT-5-mini) on SWE-bench Verified and compare text against optical input at the observation level, where the agent reads code through tool outputs. Unlike DeepSeek-OCR \cite{wei2025deepseekocr} and AgentOCR \cite{feng2026agentocr}, we require no post-training; unlike document understanding models \cite{blecher2023nougat,lee2022pix2struct}, we evaluate on a downstream software engineering task rather than transcription accuracy alone. Our study is, to our knowledge, the first to evaluate optical compression for code in an agentic software engineering setting with a full execution-based evaluation.

\begin{table}[t]
\centering
\caption{Comparison with related approaches.}
\label{tab:related-work}
\small
\begin{tabular}{p{2.0cm}cccc}
\toprule
& \textbf{Post-train} & \textbf{Code} & \textbf{Agentic} & \textbf{Zero-shot} \\
\midrule
LLMLingua \cite{jiang2023llmlingua}       & No  & No  & No  & Yes \\
DeepSeek-OCR \cite{wei2025deepseekocr}    & Yes & No  & No  & No  \\
AgentOCR \cite{feng2026agentocr}          & Yes & No  & Yes & No  \\
Nougat \cite{blecher2023nougat}           & Yes & No  & No  & No  \\
Agentless \cite{xia2024agentless}         & No  & Yes & No  & Yes \\
\textbf{Ours}                             & No  & Yes & Yes & Yes \\
\bottomrule
\end{tabular}
\end{table}
